\clearpage
\fancyhead[L]{\small DURATION OF A MAGNETIC NEEDLE’S SWING.}
\SourceLeaf{380}
\pdfbookmark[0]{Instructions for determining the duration of a magnetic needle’s swing. 1837}{schwingungsdauer-1837}
\begin{center}
{\large INSTRUCTIONS FOR DETERMINING THE DURATION OF A MAGNETIC NEEDLE’S SWING.}
\par\medskip
Resultate aus den Beobachtungen des magnetischen Vereins. 1837. IV.
\end{center}

The problem these instructions are meant to help solve is of interest in several respects. Some knowledge of the duration of a swing, even if it is not yet very precise, is already needed to apply the directions for determining the needle’s position at rest (\textit{Res.} for 1836, II.). For finding the absolute intensity of the Earth’s magnetism, however, the swing duration determined with the greatest possible precision is an essential element. In its own right, too, carrying out the operations needed to determine the duration of a swing can be regarded as good preliminary training for astronomical observations. They are especially similar to observations of stellar transits with a meridian telescope, and they have the advantage that they admit of greater precision and can be performed at any hour one chooses, without interference from unfavourable atmospheric conditions. Such observations with the magnetometer therefore also seem particularly well suited to shedding light on a subject still insufficiently understood. The results of different observers at the meridian telescope show constant differences, which can only be due to the visual or the auditory impressions, or both, reaching consciousness at different times in different people and under different circumstances.

The duration of a needle’s swing is the interval between two successive extreme positions (elongations). Strictly speaking, the needle \SourceLeaf{381}does not dwell at all at any elongation. But since the speed of its movement decreases continuously until it vanishes and then increases again just as continuously, to our senses it appears to be at rest for a longer or shorter time around the instant of elongation. With short swing durations and large arcs, however, that apparent interval of rest can scarcely be recognised as such. In every case, directly perceiving the instant of elongation in this way falls far short of the precision of an indirect determination from corresponding observations. One takes the mean of the two times at which the same division of the scale appears on the telescope’s vertical thread on the outward and return passages.

As a rule, it is best to choose a scale mark at or near the middle of the arc of the swing. The motion is fastest there, which makes the timing itself most precise. Also, when several successive swings are observed, the intervals between the recorded times are then nearly equal. In particular cases, especially when the swing duration is very long, it may sometimes be advantageous to use a different point on the scale or even several different points. We can set that possibility aside here.

For small swings it is best to choose the very division mark nearest the middle. For larger swings, prefer the nearest mark at a multiple of five or ten on the scale, since it is easier to see. For very large swings, however, the middle of the scale passes through the field of view so quickly that the chosen point on the scale has to be made sufficiently conspicuous, for example by hanging a black thread that is not too fine over the scale.

If the crossing at the thread does not coincide exactly with a one-second tick, set the fraction according to the estimated ratio of the mark’s distances from the thread at the preceding and following ticks. This is just what most astronomers are accustomed to doing when observing with a meridian telescope. It is therefore distance that one divides directly. The division of time is indirect.

Determining the swing duration from a single swing can of course give only very limited precision. It is therefore always based on a larger number of successive swings. It is true that
\expandafter\def\csname HistoricalNoteText391\endcsname{If the parts of the swinging apparatus are not firmly connected to one another, the arcs of the swings always decrease rapidly and irregularly.}
\expandafter\def\csname HistoricalNoteText382\endcsname{It hardly needs to be said that when declination is undergoing strong changes in quick succession, \textit{very small} swings (which in themselves are already ill suited for determining the duration of a swing) may show a completely distorted duration.}
\SourceLeaf{382}the swing duration depends on variable elements and is therefore, strictly speaking, itself continually changing. Apart from quite unusual cases\HistoricalNote{382}, this variability can be regarded as quite imperceptible over an interval of appreciable length. In any case, the duration of one swing deduced from a substantial number of swings will correspond to the mean values of the individual elements in question over that time. Clearly there is no need at all to follow the needle’s movements continuously over such an interval. It is enough to know the times of the first and last elongations once the swing duration is known approximately enough to leave no doubt about the \textit{number} of swings during the elapsed interval. This is further helped by the fact that one always knows in advance whether this number is even or odd, according as the first and last elongations are of the same or different kinds.

If the duration of a swing is not too short, the elongations themselves (that is, the extreme scale readings) can also be recorded between crossings, so as to derive from them the amplitudes needed for a more precise calculation of the swing duration. Moreover, the successive decrease of these amplitudes itself gives rise to some noteworthy considerations.

An example will best show how to treat the observations to obtain results from them. For this purpose, the observations made on 29 November 1835 at the astronomical observatory’s magnetometer with a twenty-five-pound bar have been chosen. The following Table I. gives the raw observations first.

\begin{center}
\begingroup\fontsize{9}{11}\selectfont\setlength{\tabcolsep}{2pt}
\begin{tabular}{r|r||r|r||r|r||r|r}
& $1755.1$ && $497.8$ && $645.9$ && $1232.1$ \\[-4pt]
$21^{\mathrm{h}}55'26''9$ && $23^{\mathrm{h}}38'49''2$ && $1^{\mathrm{h}}10'12''6$ && $2^{\mathrm{h}}49'19''7$ & \\[-4pt]
& $266.0$ && $1502.2$ && $1341.5$ && $775.9$ \\[-4pt]
$56\quad8.4$ && $39\quad31.5$ && $54.2$ && $50\quad1.5$ & \\[-4pt]
& $1751.8$ && $500.1$ && $647.3$ && $1231.0$ \\[-4pt]
$51.2$ && $40\quad13.6$ && $11\quad37.0$ && $44.1$ & \\[-4pt]
& $268.5$ && $1499.1$ && $1339.4$ && $776.4$ \\[-4pt]
$57\quad33.0$ && $56.0$ && $12\quad18.4$ && $51\quad25.8$ & \\[-4pt]
& $1748.9$ && $502.6$ && $648.7$ && $1228.7$ \\[-4pt]
$58\quad15.5$ && $41\quad38.1$ && $13\quad1.3$ && $52\quad8.5$ & \\[-4pt]
& $271.6$ && $1496.5$ && $1337.0$ && $778.0$ \\[-4pt]
$57.4$ && $42\quad20.3$ && $43.0$ && $50.0$ & \\[-4pt]
& $1744.2$ && $506.0$ && $650.7$ && $1227.0$
\end{tabular}\endgroup
\end{center}

\SourceLeaf{383}These observations, as one can see, form four sets. The first column contains the times when scale mark 1000 crossed the thread, and the second contains the elongation readings. This time they begin with the elongation preceding the first crossing and end with the elongation following the last crossing. If the elongations are omitted from the observations, it is best to note at every crossing whether the numbers passing the thread were increasing or decreasing. The method followed in Göttingen is shown here.
\begin{center}
\begin{tabular}{l}
$21^{\mathrm{h}}55'26''9-$ \\
$\phantom{21^{\mathrm{h}}}56\quad8.4+$ \\
and so forth \\
$23^{\mathrm{h}}38'49''2+$ \\
and so on
\end{tabular}
\end{center}

This is necessary to distinguish which of the elongation times deduced from the crossings correspond to minima and which to maxima. In numbering the elongation times, we have adopted the practice of marking the former with even numbers and the latter with odd numbers. The elongation time $21^{\mathrm{h}}55'47''65$, deduced from the first two crossings, is therefore preceded by the number 0, and so on.

The following Table II. gives the next results calculated from the direct observations.
\begin{center}
\begingroup\fontsize{9}{11}\selectfont\setlength{\tabcolsep}{4pt}
\begin{tabular}{r|r|r||r|r|r}
0 & $21^{\mathrm{h}}55'47''65$ & $21^{\mathrm{h}}55'47''65$ & 277 & $1^{\mathrm{h}}10'33''40$ & $21^{\mathrm{h}}55'49''54$ \\
1 & $56\quad29.80$ & $47.62$ & 278 & $11\quad15.60$ & $49.56$ \\
2 & $57\quad12.10$ & $47.74$ & 279 & $11\quad57.70$ & $49.48$ \\
3 & $57\quad54.25$ & $47.71$ & 280 & $12\quad39.85$ & $49.45$ \\
4 & $58\quad36.45$ & $47.73$ & 281 & $13\quad22.15$ & $49.57$ \\[7pt]
147 & $23\quad39\quad10.35$ & $21\quad55\quad49.89$ & 418 & $2\quad49\quad40.60$ & $21\quad55\quad49.36$ \\
148 & $39\quad52.55$ & $49.91$ & 419 & $50\quad22.80$ & $49.38$ \\
149 & $40\quad34.80$ & $49.98$ & 420 & $51\quad4.95$ & $49.35$ \\
150 & $41\quad17.05$ & $50.05$ & 421 & $51\quad47.15$ & $49.37$ \\
151 & $41\quad59.20$ & $50.02$ & 422 & $52\quad29.25$ & $49.29$
\end{tabular}\endgroup
\end{center}

The second column gives the resulting elongation times. The numbers in the first column follow directly for the first five elongations. For the later ones they are found as follows.

\SourceLeaf{384}Comparing the. elongation 0 with 4 gives $42''20$ as an approximate swing duration. If we divide the interval between elongation 4 and the next one, $1^{\mathrm{h}}40'33''90$, by this value, and recall that the index of the latter must be odd, the quotient $142.983$ leaves no doubt that 143 swings must have elapsed between those two elongations. Indeed, if we assumed 141 or 145, the resulting swing duration would be $42''7936$ or $41''6131$, both too far from the approximate value $42''20$ to be admissible. Starting from 143 swings gives a duration of $42''1951$. This could be used as the basis for numbering the subsequent sets of observations, although in the present case, where there are no \textit{very} long interruptions, the first approximation already suffices throughout.

To obtain the swing duration more accurately, and even to detect its variation over the whole series of observations, one can first compare the first set with the second as follows. The duration of 147 swings is obtained from
\begin{center}
\begin{tabular}{rcl}
$0\text{—}147$ & $.\ .\ .\ .\ .$ & $1^{\mathrm{h}}43'22''70$\\
$1\text{—}148$ & & $\phantom{1^{\mathrm{h}}43'}22.75$\\
$2\text{—}149$ & & $\phantom{1^{\mathrm{h}}43'}22.70$\\
$3\text{—}150$ & & $\phantom{1^{\mathrm{h}}43'}22.80$\\
$4\text{—}151$ & & $\phantom{1^{\mathrm{h}}43'}22.75$
\end{tabular}
\end{center}
giving a mean of $1^{\mathrm{h}}43'22''74$, or the duration of one swing $=42''19551$. In the same way, the swing duration between the second and third sets is $=42''17654$, between the third and fourth $=42''17879$, and between the first and fourth, or the mean over the whole series, $=42''18344$.

This calculation can also be carried out in a slightly different form, with the added advantage of showing clearly the regular course of all the individual observations. One starts by reducing the individual elongation times to a common epoch using an approximate swing duration, calculating back from each time by the amount taken up by all the swings since that epoch, with the aid of this approximate value. One therefore subtracts from each number in the second
\SourceLeaf{385}column the product of this assumed value and the corresponding number in the first column. If the observations were absolutely accurate, and the swing duration exactly constant and exactly equal to the assumed value, all the numbers reduced in this way would be exactly equal. An increase in the numbers from one set to the next, on the other hand, shows that the swing duration used as the basis for that interval was too small, and conversely. The irregular fluctuations of the results within one and the same set give a measure of the accuracy of the observations themselves.

In our example, comparing the first elongation time with the last gives a swing duration $=42''18389$. The approximate value $42''18$ was used in its place to calculate the numbers in the third column. One can see at a glance that this swing duration is slightly too small for the interval from the first to the second set, but a little too large from the second to the third. and likewise from the third to the fourth. For accurate results, take the mean of the third-column numbers belonging to each set. These means
\begin{center}
\begin{tabular}{r}
$21^{\mathrm{h}}55'47''69$\\
$49.97$\\
$49.52$\\
$49.35$
\end{tabular}
\end{center}
can be regarded as more accurate values of the reduced times corresponding to the indices 2, 149, 279, 420. Thus from the first set to the second the observations run $2''28$ ahead of the assumed course over 147 swings, which amounts to $0''01551$ for one swing. The corrected value is therefore $42''19551$, in exact agreement with the value found above. Comparing the subsequent sets gives the same result.

The quality of the observations themselves is evident simply from looking at the numbers belonging to one set. It is nevertheless worth giving directions here for finding a more definite measure of accuracy in suitable cases. If we denote the mean error to be expected at a crossing by $\varepsilon$, the number of results belonging to a set by $s$, and the sum of the squares of the differences of those individual results
\SourceLeaf{386}from their mean by $q$, we may assume approximately
\[
\frac{(s-1)^2\varepsilon\varepsilon}{2s}=q
\]
or, if there are several sets,
\[
\varepsilon\varepsilon\sum\frac{(s-1)^2}{2s}=\sum q
\]
hence
\[
\varepsilon=\BareRoot{\frac{\sum q}{\sum\frac{(s-1)^2}{2s}}}
\]

In our example the differences from the mean for the first set are $0''04$, $0''07$, $0''05$, $0''02$, $0''04$, so $q=110$ if one takes a hundredth of a second as the unit. Further, $s=5$, hence $\frac{(s-1)^2}{2s}=\frac85$. For the three subsequent sets, with the same value of $s$, we have $q=190;\ 110;\ 50$. We therefore have
\[
\tfrac{32}{5}\varepsilon\varepsilon=460
\]
or
\[
\varepsilon=8.5,\qquad\text{that is, }\varepsilon=0''085.
\]

It must be noted, however, that the validity of this rule depends on several conditions that our example does not sufficiently meet. First, the approximate swing duration assumed in calculating the reduced numbers must be one that can be regarded, without appreciable error, as the true duration during each set. Second, the different sets that are combined must have been observed under nearly the same conditions, insofar as those conditions can affect the accuracy of observation. Neither holds in our example, so the calculation above should be regarded only as an illustration of the formula. For more accurate determinations, it is better first to make observations specifically for that purpose. Leaving a fuller treatment of this subject for the future, let it be noted here only that the accuracy of observation varies with the particular instrument and observer, the quality of the lighting on the scale, the speed of the swing and the kind of clock. Both a movement that is too fast and one that is too slow are less favourable to the accuracy of
\SourceLeaf{387}observation than a moderate speed. Observations made with a clock that ticks every second are less precise than those made with a chronometer that marks shorter intervals. In the most favourable circumstances the accuracy of these observations far exceeds that of the best observations with a meridian telescope.

As is well known, all else being equal, the smaller the arc of the swing, the shorter the swing duration. As the arc approaches zero without limit, the duration has a limiting value. If we denote this limiting value, or in the usual language the duration of an infinitesimally small swing, by $T$, and the duration corresponding to an arc $G$ by $T'$, we have the familiar expression
\[
T'=T(1+\tfrac14\sin\tfrac14G^2+\tfrac14.\tfrac9{16}\sin\tfrac14G^4+\tfrac14.\tfrac9{16}.\tfrac{25}{36}\sin\tfrac14G^6+\text{etc.})
\]

Given the small arcs to which one is confined when using the magnetometer, one can safely set aside terms of the fourth and higher orders. One can therefore also put $\frac{g}{8r}$ in place of $\sin\tfrac14G$, where $g$ is the part of the scale corresponding to the arc $G$, and $r$ is the horizontal distance from the midpoint of the scale to the mirror. We thus have $T'=T(1+\frac{gg}{256rr})$, or, to the same accuracy, $T=T'(1-\frac{gg}{256rr})$. In our example, measured in scale divisions or millimetres, $r=4775.9$. The arc of the swing between elongations 0 and 1 is $=1485.8$, so the reduction of the swing duration to that of an infinitesimally small swing is $=-0''01595$. The reductions for the second, third and fourth swings are similarly found to be $0''01590$, $0''01584$ and $0''01577$, respectively. The mean of the first four swing durations reduced to infinitesimally small swings is thus $=42''18414$.

This method, however, requires that the elongations be observed without interruption. When only the beginning and end of a series of swings have been observed, the reduction to infinitesimally small arcs may, if the arc has decreased only moderately in the meantime, be carried out if need be by taking a mean value of its magnitude. But reducing a longer series of such swings, in which the arc has decreased considerably, necessarily requires at least an approximately correct
\SourceLeaf{388}knowledge of the law governing the decrease. I therefore turn to the treatment of the observed elongations, the first results of which are given in Table III below.
\begin{center}
\fontsize{9}{11}\selectfont
$\left.\begin{array}{r|r|r|r}
0&1009.725&1487.45&3.17244\\
1&1009.525&1484.55&3.17160\\
2&1009.425&1481.85&3.17081\\
3&1009.475&1478.85&3.16992\\
4&1009.075&1474.95&3.16878
\end{array}\right\}3.170710$\\[6pt]
$\left.\begin{array}{r|r|r|r}
147&1000.575&1003.25&3.00141\\
148&1000.375&1000.55&3.00024\\
149&1000.225&997.75&2.99902\\
150&1000.200&995.20&2.99791\\
151&1000.400&992.20&2.99660
\end{array}\right\}2.999036$\\[6pt]
$\left.\begin{array}{r|r|r|r}
277&994.050&694.90&2.84192\\
278&993.875&693.15&2.84083\\
279&993.700&691.40&2.83973\\
280&993.450&689.50&2.83853\\
281&993.350&687.30&2.83715
\end{array}\right\}3.839630$\\[6pt]
$\left.\begin{array}{r|r|r|r}
418&1003.725&455.65&2.65863\\
419&1003.575&454.85&2.65787\\
420&1003.125&453.45&2.65653\\
421&1002.950&451.50&2.65466\\
422&1002.925&449.85&2.65307
\end{array}\right\}2.656152$
\end{center}

The first column contains the index of each elongation. The second gives the corresponding position of the needle at rest, by the formula $\tfrac14(a+2b+c)$, where $b$ is the observed elongation and $a$ and $c$ are the preceding and following ones (see \textit{Resultate} for 1836, II). The third column gives twice the distance of each elongation from the corresponding position at rest, or the arc that would have been described from that point in the absence of the causes tending to reduce it. This is $\tfrac12(a-2b+c)$ or $\tfrac12(2b-a-c)$, that is, the mean of the preceding and following arcs. The fourth column contains the logarithm of this number. Finally, alongside it is the mean of the fourth-column numbers belonging to each set.

All experience agrees in indicating that, at least over a moderately long time, one may regard the numbers in the third column as decreasing in geometric progression,
\SourceLeaf{389}and hence their logarithms as decreasing in arithmetic progression, or at least accept this as the most plausible approximation. The small irregularities found in comparing successive numbers within a set are attributable only to unavoidable small errors of observation or small chance disturbances. Their detrimental influence is reduced as far as possible by treating the means in the fifth column as belonging to the corresponding middle indices, and then deducing from them the course over the whole series of observations.

We accordingly have the following logarithms of the amplitudes for the elongations
\begin{center}
\begin{tabular}{rcl}
$2$ & $.\ .\ .\ .\ .\ .\ .\ .\ .\ .$ & $3.170710$\\
$149$ & $.\ .\ .\ .\ .\ .\ .\ .\ .\ .$ & $2.999036$\\
$279$ & $.\ .\ .\ .\ .\ .\ .\ .\ .\ .$ & $2.839630$\\
$420$ & $.\ .\ .\ .\ .\ .\ .\ .\ .\ .$ & $2.656152$
\end{tabular}
\end{center}

The logarithm therefore decreased over 147 swings from the first set to the second by $0.171674$. Spread uniformly, that amounts to $0.00116785$ per swing. I call this quotient the \textit{logarithmic decrement}. From the second to the third set it is found to be $=0.00122620$, and from the third to the fourth $=0.00130126$. At least in this case the decrease is not far from uniform. I shall return below to the variability of the logarithmic decrement.

Now, on the assumption that the amplitudes over a series of swings have decreased in geometric progression, it is easy to reduce the swings to infinitesimally small ones. Let $g$ be the magnitude of the first swing in scale divisions, $g^0$ the magnitude of the last, and $\theta$ the exponent of the geometric progression, so that $g^0=g\theta^\mu$, where $\mu$ is the number of swings. Then the reduction of the first swing duration to that of an infinitesimally small swing is
\[
-\frac{Tgg}{256rr},
\]
that of the second is
\[
=-\frac{Tgg\theta\theta}{256rr}
\]
and so on. The sum of them all is therefore
\[
=-\frac{T(gg-g^0g^0\theta\theta)}{256rr(1-\theta\theta)}
\]

\SourceLeaf{390}If we denote by $h$ and $h^0$ the amplitudes corresponding to the initial and final elongations, calculated in the same way as in Table III, then $h=\tfrac12(\frac{g}{\theta}+g)$, $h^0=\tfrac12(g^0+g^0\theta)$, and hence the sum above is
\[
=-\frac{T(hh-h^0h^0)}{64rr}.\frac{\theta\theta}{(1+\theta)^2(1-\theta\theta)}
\]

In all the cases under consideration, $\theta$ differs little from unity, so the logarithmic decrement, for which we write $\lambda$, is a small fraction, $=\log\frac1\theta$. If $m$ denotes the modulus of the system of logarithms, the second factor in that expression has the following series expansion
\[
\frac{\theta\theta}{(1+\theta)^2(1-\theta\theta)}=\frac{m}{8\lambda}-\tfrac5{96}.\frac{\lambda}{m}+\tfrac{37}{2880}(\frac{\lambda}{m})^3\ \text{etc.}
\]
For sufficient accuracy we may replace that factor by just its first term $\frac{m}{8\lambda}$. The reduction of the total duration of the $\mu$ swings is therefore
\[
=-\frac{Tm(hh-h^0h^0)}{512rr\lambda}
\]
and the reduction of the average for one swing is
\[
=-\frac{Tm(hh-h^0h^0)}{512rr\lambda\mu}
\]

For the reasons given above, it makes no difference in these formulas whether we use the corrected or the uncorrected value of $T$. In our example we obtain the following values for the reduction
\begin{center}
\fontsize{9}{11}\selectfont
\begin{tabular}{r|r|r|r}
& \multicolumn{1}{c|}{For the full} & \multicolumn{1}{c|}{For one} & \multicolumn{1}{c}{Reduced}\\
& \multicolumn{1}{c|}{interval} & \multicolumn{1}{c|}{swing} & \multicolumn{1}{c}{value}\\[5pt]
$2\ .\ .\ .\ .\ 149$ & $-1''6111$ & $-0''01096$ & $42''18455$\\
$149\ .\ .\ .\ .\ 279$ & $-0.6624$ & $-0.00510$ & $42.17144$\\
$279\ .\ .\ .\ .\ 420$ & $-0.3285$ & $-0.00233$ & $42.17646$
\end{tabular}
\end{center}

The decrease in the arc of swing is of interest in several other respects besides its bearing on a more accurate calculation of the swing duration. We must first distinguish the external circumstances in which the needle swings.

\SourceLeaf{391}If the apparatus is properly set up and in perfectly good condition\HistoricalNote{391}, and there is nothing in its immediate surroundings that must appreciably damp its swinging motion, the arcs of swing always decrease very slowly. The decrease is regular in the sense that, for a moderate length of time at least, it follows a geometric progression, so the logarithmic decrement is nearly constant. In our example this logarithmic decrement changes over nearly six hours only from $0.00117$ to $0.00130$, or the decrease in the arc from one swing to the next varied from $\tfrac1{371}$ to $\tfrac1{334}$. Yet experience shows that very different values of the logarithmic decrement occur very often. In the same apparatus it can rise to around $0.00300$, and at other times it falls to $0.00030$ and, in rare cases, even lower. In our experience, however, the changes are always gradual. This decrease can therefore hardly be attributed to air resistance alone. The actual cause of these variations has so far eluded our many repeated experiments, and we therefore very much hope that observers elsewhere will give particular attention to this still puzzling phenomenon. We have noticed one circumstance so often in these experiments that we can hardly still regard it as accidental, although any causal connection remains entirely unexplained. The \textit{very} small values of the logarithmic decrement always occur only when the sky is overcast, whereas the very large ones generally occur only when it is clear. We may add, for good measure, that the apparatus is suspended from a metal wire rather than a silk thread, and that these experiments were carried out only on windless days. Both hygrometric effects and drafts are therefore out of the question in these two respects.

If, on the other hand, the needle is surrounded by a multiplier that forms part of a closed circuit, there is an additional cause for the decrease in the arcs of swing. The motion of the needle induces a galvanic current in the circuit. The current is strongest when the
\expandafter\def\csname HistoricalNoteText393\endcsname{Whatever means are used for this, one cannot expect to produce strictly \textit{pure} swings about a vertical axis. For this reason alone, one should not expect the results to agree with the utmost precision, which in any case is unnecessary here.}
\expandafter\def\csname HistoricalNoteText392\endcsname{It is not exactly the same in all experiments, since the effect of the multiplier combines with a quantity that is not itself invariable, as noted above. The effect of the multiplier itself also depends in part on the conductivity of the wire, which varies with temperature.}
\SourceLeaf{392}swinging velocity of the needle is greatest. It reverses direction as soon as the needle does. The reaction of this current on the needle, however, always consists in a reduction of the needle’s swinging velocity, and theory shows that this too must cause the arc of swing to decrease in a very nearly geometric progression. Since the intensity of the induced current also depends on the resistance offered by the \textit{whole} circuit, the increase in the logarithmic decrement produced by the multiplier is greatest when the circuit is closed immediately beyond it. The larger the part added to the circuit, the smaller the increase. With an open circuit there is no effect at all, and the logarithmic decrement is the same as if the multiplier had been removed altogether. The apparatus used for the observations above has a multiplier with 610 turns. If its circuit is closed without adding anything further, the logarithmic decrement rises to about $0.02400$\HistoricalNote{392}, so that the arc of swing is reduced by half in only about 9 minutes. With the circuit open and with a value of the logarithmic decrement such as the observations above gave, this takes about three hours. In the case of the four-pound needle at the magnetic observatory, closing the circuit of a multiplier with 536 turns produces a logarithmic decrement of nearly the same magnitude. Since that needle has a swing duration of $21''6$, however, the arc of swing is already halved after $4\tfrac12$ minutes. With such strong damping, once the needle has settled, no appreciable swings can develop unless there are extraordinary disturbances from outside or unusually large and rapid changes in declination. Besides its countless other uses, such a powerful multiplier therefore offers the important advantage of greatly facilitating the scheduled observations and making all other means of bringing the needle to rest unnecessary.

The device mentioned above [p. 372 of this volume] under the name of a damper produces this effect to a much greater degree. The damper made for the magnetome\SourceLeaf{393}ter of the magnetic observatory consists of two oblong rectangular copper frames, each weighing 13 pounds. They can be slid like a lining into the wooden frames of the two halves of the multiplier. This device proves extremely useful in very many cases, and a similar accessory has already been supplied with several magnetometers delivered by Mr. \textsc{Meyerstein} to observers elsewhere. A few remarks on it will therefore be in place here.

First of all, anyone who intends to use such an apparatus will want to know its damping strength in numbers. The needle is first set swinging through very large arcs\HistoricalNote{393}. The elongations are recorded as soon as they fall within the scale and for as long as the arc of swing still has an appreciable magnitude. If $g$, $g'$, $g''$, $g'''$ and so on are the resulting amplitudes in scale divisions, the differences of their logarithms $\log g-\log g'$, $\log g'-\log g''$, $\log g''-\log g'''$ and so on give just as many determinations of the logarithmic decrement. When the declination is changing appreciably, this change is eliminated, on the assumption that it is uniform during the observations, by using the formulas
\[
\log(g+g')-\log(g'+g''),\quad \log(g'+g'')-\log(g''+g''')\text{ and so on}
\]
As is easy to see, this comes to the same thing as the procedure given above [p. 383 of this volume].

One of the experiments carried out on January 9, 1838, gave, for example, the following results
\begin{center}
\begin{tabular}{r|r|r|r}
Elongations&Amplitudes&Logarithms&Differences\\
$134.0$&&&\\[-3pt]
&$1057.7$&$3.02436$&\\[-3pt]
$1191.7$&&&$0.31137$\\[-3pt]
&$516.4$&$2.71299$&\\[-3pt]
$675.3$&&&$0.30850$\\[-3pt]
&$253.8$&$2.40449$&\\[-3pt]
$929.1$&&&$0.31141$\\[-3pt]
&$123.9$&$2.09307$&\\[-3pt]
$805.2$&&&
\end{tabular}
\end{center}
\SourceLeaf{394}Thus the mean logarithmic decrement was $=0.31043$. The other formulas give
\[
\log\frac{1574.1}{770.2}=0.31043,\quad \log\frac{770.2}{377.7}=0.30945
\]
and hence a mean of $0.30994$. In these experiments the multiplier circuit was open, so the damper alone was acting. Other experiments on the same day, in which the multiplier circuit was closed as well, gave a logarithmic decrement $=0.33570$.

Since copper, when it is not entirely pure, can exert a direct magnetic influence, however weak, it is inadvisable to use the damper for observations aimed at determining the absolute declination before making sure that there is no appreciable influence of this kind. To find out, one observes the position of the needle alternately with and without the damper while simultaneously observing the position of a second needle at a suitable distance, so that changes in declination during the observations can be accounted for. In the absence of a second magnetometer, one can eliminate these changes, though less reliably, by taking the alternating readings at nearly equal intervals and comparing each position without the damper with the mean of the preceding and following positions with the damper, and vice versa. In making the damper for the magnetic observatory here, Mr. \textsc{Meyerstein} took the precaution of using only completely new files. The finished damper was then left for a considerable time in dilute hydrochloric acid, then put in an alkaline solution and finally rinsed with water. Experiments of the kind described showed no appreciable influence of this damper on the position of the needle.

What would hold strictly if the arc of swing always stayed the same undergoes several modifications in the actual case, where the arc of swing continually decreases. These deserve a somewhat closer look here, if only to determine precisely under what circumstances they may be regarded as imperceptible.

A decrease in the arc of swing in geometric progression presupposes a force opposing the motion at each instant, in simple proportion to
\expandafter\def\csname HistoricalNoteText395\endcsname{\textit{Strictly} speaking, both hold only for infinitesimally small swings.}%
\SourceLeaf{395}its velocity\HistoricalNote{395}. The general equation for the swinging motion therefore has the following form if we neglect quantities of the third order with respect to the arc of swing
\[
0=\frac{\mathrm{d}\mathrm{d}x}{\mathrm{d}t^2}+nn(x-p)+2\epsilon.\frac{\mathrm{d}x}{\mathrm{d}t}
\]
Here $x$ denotes the scale division marking the position of the needle at time $t$, and $p$ the scale division corresponding to the position of rest. The quantities $nn$ and $2\epsilon$ denote the magnetic directive force and the retarding force, both divided by the needle’s moment of inertia. The complete integral of this equation is
\[
x=p+Ae^{-\epsilon t}\sin\BareRoot{(nn-\epsilon\epsilon)}.(t-B)
\]
where $e$ is the base of the natural logarithms and $A$ and $B$ are the two arbitrary constants introduced by integration. Without the retarding force the integral would be
\[
x=p+A\sin n(t-B)
\]
Thus the needle oscillates periodically about the point $p$ in the former case as in the latter, but there is a twofold difference. In the second case the greatest deviation from the centre, or half the amplitude, is constant $=A$, whereas in the first it decreases in geometric progression. And in the first case the argument of the periodic function increases more slowly than in the second. If we set the swing duration in the second case, where it depends only on the magnetic directive force, $=T$, and in the first $=T'$, then, taking $\pi$ in its usual sense, we have
\[
nT=\pi,\quad \BareRoot{(nn-\epsilon\epsilon)}.T'=\pi
\]
If for brevity we write $n'$ in place of $\BareRoot{(nn-\epsilon\epsilon)}$ and introduce an auxiliary angle $\varphi$ so that
\[
\sin\varphi=\frac{\epsilon}{n},\quad \cos\varphi=\frac{n'}{n},\quad \operatorname{tang}\varphi=\frac{\epsilon}{n'}
\]
and if, as above, we denote the logarithmic decrement by $\lambda$ and the modulus of the system by
\SourceLeaf{396}$m$, we obtain
\[
\begin{gathered}
T'=\frac{T}{\cos\varphi}\\
\frac{\lambda}{m}=\epsilon T'=n'\operatorname{tang}\varphi.T'=\pi\operatorname{tang}\varphi
\end{gathered}
\]
and consequently
\[
\operatorname{tang}\varphi=\frac{\lambda}{m\pi}=\frac{\lambda}{1.364376}
\]
For $\lambda=0.02400$ and $T=42''18$ these formulas give $\varphi=1^0\,0'\,28''$ and $T'=42''18653$. Merely closing the multiplier circuit therefore produces only a slight increase in the swing duration. On the other hand, if we set $T=20''60$, the figures found above when the damper was used alone or together with the multiplier give
\begin{center}
\begin{tabular}{r|r|r}
$\lambda=0.30994$&$\varphi=12^0\,47'\,54''$&$T'=21''12484$\\
$\lambda=0.33570$&$\varphi=13\quad49\quad22$&$T'=21.21439$
\end{tabular}
\end{center}

The observations agree with this calculated increase in swing duration as closely as one could expect from the small number of swings to which one must confine oneself here.

When the arcs of swing decrease, the true times of the elongations differ slightly from the times derived from corresponding positions. With logarithmic decrements as large as those found when a damper is used, this difference becomes fairly substantial.

In the integral given above, $B$ evidently denotes the time of a passage through the position of rest $p$. Since the choice of the instant from which time is reckoned makes no difference, we shall set $B=0$ for greater simplicity. Our formula then becomes
\[
x=p+Ae^{-\epsilon t}\sin n't
\]
The next passage through $p$ after the one at $t=0$ takes place at $n't=180^0$, or $t=T'$. The time of the elongation obtained simply from these corresponding observations is therefore $t=\tfrac12T'$, whereas the needle actually comes to a stop earlier. For $\frac{\mathrm{d}x}{\mathrm{d}t}=0$ we have
\[
0=Ae^{-\epsilon t}(-\epsilon\sin n't+n'\cos n't)
\]
\SourceLeaf{397}Hence
\[
\operatorname{cotang}n't=\frac{\epsilon}{n'}=\operatorname{tang}\varphi
\]
The first positive value is therefore $n't=\tfrac12\pi-\varphi$, and $t=\tfrac12T'-\frac{\varphi T'}{\pi}$, or, with $\varphi$ expressed in degrees,
\[
t=\tfrac12T'-\frac{\varphi}{180^0}.T'
\]

The same difference evidently occurs at the next elongation. From the figures given above it is found to be $=0''23$ for the twenty-five-pound bar with the multiplier in use, $=1''50$ for the magnetometer of the M. O. with the damper alone, and $=1''63$ when the damper and multiplier are used together.

The actual stops occur earlier than the instants inferred from successive passages through the position of rest $p$ by a constant time independent of the size of the arc of swing. We may therefore simply retain the latter instants, since only the \textit{differences} between the elongation times matter when determining the swing duration. We must take care, however, to choose the point $p$ itself or a point very close to it for observation, and have all the more reason to exclude small arcs of swing. With such arcs and large logarithmic decrements even a slight departure from the correct $p$ would produce an appreciable error. Though that condition is easy to meet, it may still be worth giving the general formula here for the error in the elongation times inferred from corresponding passages.

Let $u$ be half the interval between two corresponding passages through a point $x$, and let $\tfrac12T'-\delta$ be the mean of the passage times, or the elongation time inferred from them. The passage times themselves are then $\tfrac12T'-u-\delta$ and $\tfrac12T'+u-\delta$, so we have
\[
\begin{aligned}
x&=p+Ae^{-\epsilon(\tfrac12T'-u-\delta)}\sin n'(\tfrac12T'-u-\delta)\\
x&=p+Ae^{-\epsilon(\tfrac12T'+u-\delta)}\sin n'(\tfrac12T'+u-\delta)
\end{aligned}
\]
from which it follows, since $\tfrac12n'T'=\tfrac12\pi$, that
\[
e^{2\epsilon u}\cos n'(u+\delta)=\cos n'(u-\delta)
\]
\SourceLeaf{398}and hence
\[
\operatorname{tang}n'\delta=\frac{e^{2\epsilon u}-1}{(e^{2\epsilon u}+1)\operatorname{tang}n'u}
\]

For $u=\tfrac12T'$ this formula gives $\delta=0$. This is the case in which the position of rest $p$ itself is chosen for the passages. An infinitesimally small value, on the other hand, corresponds to the true stopping point, and here the formula gives
\[
\operatorname{tang}n'\delta=\frac{\epsilon}{n'}=\operatorname{tang}\varphi
\]
Thus $\delta=\frac{\varphi}{n'}=\frac{\varphi T'}{\pi}$, in agreement with what was found above.

Finally, when the arcs of swing decrease, the calculation of the observations relating to the needle’s position of rest also needs a modification. It becomes appreciable, of course, only when the swings diminish very substantially.

The positions of the needle $x$, $x'$, corresponding to two times $t$, $t+T'$ separated by one swing duration, have the values
\[
\begin{aligned}
x&=p+Ae^{-\epsilon t}\sin n't\\
x'&=p+Ae^{-\epsilon t-\epsilon T'}\sin(n't+n'T')
\end{aligned}
\]
or, since $n'T'=\pi$,
\[
x'=p-Ae^{-\epsilon t-\epsilon T'}\sin n't
\]
or, if as above we denote by $\theta$ the fraction whose Briggs logarithm is $-\lambda$ and whose natural logarithm is therefore $\epsilon T'$,
\[
x'=p-\theta Ae^{-\epsilon t}\sin n't
\]

It is clear, then, that to find $p$ we can no longer take the arithmetic mean of $x$ and $x'$. We must divide the difference between $x$ and $x$ in the ratio of 1 to $\theta$. This gives
\[
p=\frac{\theta x+x'}{1+\theta}=x+\frac{1}{1+\theta}.(x'-x)=x'-\frac{\theta}{1+\theta}.(x'-x)
\]
Since in these observations the difference $x'-x$ is always very small, for ease of calculation we may allow ourselves to use in place of $\frac{\theta}{1+\theta}$ a nearby fraction expressed in small numbers.
\SourceLeaf{399}For example, for $\lambda=0.30994$ we may use, in place of the more accurate value $0.3288$, the approximation $\tfrac13$.

The question now arises at what instant this result is to be regarded as valid. When the arcs of swing decrease only very slowly, the simple mean of the scale readings is taken to correspond to the simple mean of the times. It might likewise seem that when the difference between the positions is divided unequally, the time interval should be divided in the same ratio, so that the position $x+\frac{x'-x}{1+\theta}$ would be valid at $t+\frac{T'}{1+\theta}$. This is theoretically incorrect, however, and a closer examination, though of no practical consequence in ordinary cases, may still deserve a place here on theoretical grounds.

Evidently the instant at which the result holds matters only insofar as we are unwilling to regard the magnetic declination as constant during the interval between the two records. We can always regard it as changing uniformly during that interval, however, and must do so if the calculation is to be given a definite basis\quad In this case our fundamental equation thus has the form
\[
0=\frac{\mathrm{d}\mathrm{d}x}{\mathrm{d}t^2}+nn(x-p-\alpha t)+2\epsilon.\frac{\mathrm{d}x}{\mathrm{d}t}
\]
whose complete integral is
\[
x=p-\frac{2\alpha\epsilon}{nn}+\alpha t+Ae^{-\epsilon t}\sin n'(t-B)
\]
if, as above, $n'$ stands for $\BareRoot{(nn-\epsilon\epsilon)}$. If $x'$ then denotes the position at time $t+T'$, we obtain, with $\theta$ in its previous sense,
\[
x'=p-\frac{2\alpha\epsilon}{nn}+\alpha t+\alpha T'-\theta Ae^{-\epsilon t}\sin n'(t-B)
\]
and consequently
\[
\frac{\theta x+x'}{1+\theta}=p-\frac{2\alpha\epsilon}{nn}+\alpha t+\frac{\alpha T'}{1+\theta}
\]
This result is therefore the position of rest at the time
\[
t+\frac{T'}{1+\theta}-\frac{2\epsilon}{nn}=t+\frac{T'}{1+\theta}-\frac{\sin2\varphi.T'}{\pi}
\]
Thus for $\lambda=0.30994$ the time for which the result calculated by the above rule holds is $=t+0.5337T'$, and for $\lambda=0.33570$ it is $=t+0.5365T'$.
\SourceLeaf{400}We see, then, that even with such strong damping the instant for which the result holds differs only slightly from the simple mean of the times.

Everything developed so far rests on the assumption that $\epsilon$ is smaller than $n$. In the opposite case the integral of the fundamental equation takes another form. In place of the term $Ae^{-\epsilon t}\sin\BareRoot{(nn-\epsilon\epsilon)}.(t-B)$, when $\epsilon$ is greater than $n$, we obtain two terms of the form
\[
Ae^{-(\epsilon+\BareRoot{(\epsilon\epsilon-nn)})t}+Be^{-(\epsilon-\BareRoot{(\epsilon\epsilon-nn)})t}
\]
and when $\epsilon=n$, a term of the following form
\[
(A+Bt)e^{-\epsilon t}
\]

In either case the motion has nothing periodic about it, and the position approaches the position of rest asymptotically. For our damper $\frac{\epsilon}{n}=0.22152$, so a damper acting more than $4\tfrac12$ times as strongly would have to be used to produce that effect. Merely increasing the amount of metal in the same proportion would evidently be insufficient, since the extra metal would have to be added on the outside and the outer layers of the metal frame contribute relatively less to the inductive effect than the inner layers. It would, in fact, be inadvisable to use damping strong enough for the motion to cease to be periodic. For one thing, as soon as $\epsilon$ exceeds the limiting value $n$, the approach to the position of rest becomes slower again. For another, we would lose the essential advantage of being able to calculate the position of rest conveniently from any two records separated by $T'$.

